\documentclass{article}
\usepackage[T1]{fontenc}
\usepackage{lmodern}
\usepackage{graphicx}
\usepackage{amsmath,amssymb}
\usepackage[a4paper,margin=2cm]{geometry}
\usepackage[absolute,overlay]{textpos}
\setlength{\TPHorizModule}{1mm}
\setlength{\TPVertModule}{1mm}
\usepackage[indonesian]{babel}
\usepackage{hyperref}
\setlength{\emergencystretch}{2em}
% unit: Halaman 27

\begin{document}

% === Halaman 27 (python/native) ===

(02 • 26) = (0000 0010)  (0010 0110) 

                = x  (x5 + x2 + x)  mod (x8 + x4 + x3 + x + 1)


   = (x6 + x3 + x2) mod (x8 + x4 + x3 + x + 1)


   = x6 + x3 + x2


   = (01001100)


   = 4C

(03 • 7B) = (0000 0011)  (0111 1011) 

                 = (x + 1)  (x6 + x5 + x4 + x3 + x + 1) mod (x8 + x4  +  x3 + x + 1)



    = ((x7 + x6 + x5 + x4 + x2 + x) + (x6 + x5 + x4 + x3 + x  +1)) mod (x8 + x4 +  x3 + x + 1)

                = (x7 + (1 + 1)x6 + (1 + 1)x5 + (1 + 1)x4 + x3 + x2 + (1 +1)x + 1) mod (x8 + x4 + x3 + 

                    x + 1)

                = (x7 + x3 + x2 + 1) mod (x8 + x4 + x3 + x + 1)


   = (x7 + x3 + x2 + 1)


   = (1000 1101)  = 8D    

(01 • BD) = BD = 10111101

(01 • 43) = 43 = 01000011

(02 • 26)  (03 • 7B)  (01 • BD)  (01 • 43) = 3F

\end{document}
